\newcommand{\figuraSeparacionMolecular}{%
\begin{figure}[htbp]
    \centering
    \begin{tikzpicture}[>=Latex, font=\small]
        \draw[thick, rounded corners=2pt] (0,0) rectangle (6,4.5);
        \node[font=\bfseries] at (3,4.85) {Volumen macroscópico $V$};

        \foreach \x/\y in {
            0.55/0.55,1.50/0.80,2.55/0.55,3.50/0.85,4.55/0.55,5.45/0.90,
            0.80/1.60,1.80/1.75,2.80/1.55,3.85/1.75,4.85/1.55,5.55/1.90,
            0.50/2.60,1.45/2.75,2.45/2.55,3.75/2.55,4.70/2.75,5.45/2.55,
            0.80/3.65,1.85/3.55,2.85/3.75,3.90/3.55,4.95/3.70,5.55/3.45}
        {
            \fill[blue!65] (\x,\y) circle (0.11);
        }

        \fill[red!75] (2.45,2.55) circle (0.14);
        \fill[red!75] (3.75,2.55) circle (0.14);
        \draw[<->,thick,red!70!black]
            (2.59,2.55) -- node[above,fill=white,inner sep=1pt]
            {$L_{\mathrm{mol}}$} (3.61,2.55);
        \node at (3,0.20) {$N$ moléculas};

        \coordinate (A) at (7.25,0.65);
        \coordinate (B) at (10.05,0.65);
        \coordinate (C) at (10.05,3.45);
        \coordinate (D) at (7.25,3.45);
        \coordinate (E) at (7.85,1.25);
        \coordinate (F) at (10.65,1.25);
        \coordinate (G) at (10.65,4.05);
        \coordinate (H) at (7.85,4.05);

        \draw[thick] (A)--(B)--(C)--(D)--cycle;
        \draw[thick] (B)--(F)--(G)--(C);
        \draw[thick] (D)--(H)--(G);
        \draw[dashed] (A)--(E)--(F);
        \draw[dashed] (E)--(H);

        \shade[ball color=red!70] (8.95,2.35) circle (0.22);
        \node[font=\bfseries] at (8.95,4.65) {Celda molecular promedio};
        \draw[<->,thick]
            (7.25,0.30)--node[below]{$L_{\mathrm{mol}}$}(10.05,0.30);
        \node[align=center] at (8.95,-0.55)
            {$V_{\mathrm{mol}}\approx L_{\mathrm{mol}}^3$\\[2pt]
             $V_{\mathrm{mol}}=V/N$};
    \end{tikzpicture}
    \caption{
    Cada molécula se asocia con una celda de volumen promedio $V/N$; al
    idealizarla como un cubo, su lado corresponde a $L_{\mathrm{mol}}$.}
    \label{fig:separacion-molecular}
\end{figure}%
}